\chapter{“市民社会”概念及其历史起源问题}

如果比较一下1820年至1850年间黑格尔同时代人对（当时即已引发不少赞叹的）《法哲学原理》第三部分所阐述的市民社会的评论与当今解释者的评论，那么对这两拨人来说，这个概念的历史地位似乎不是一个问题。对于像甘斯这样的黑格尔的亲炙弟子来说（甘斯1833年在“逝者友人版”《黑格尔全集》总体框架下编辑出版了《法哲学原理》第2版），不言而喻的是，这里“与国家有关的一切，包举无遗”，因为“甚至国民经济的科学”都能“在市民社会中找到恰如其分的地位和讨论”。\footnote{黑格尔：《法哲学原理》，序言，引自《全集》第8卷(1833年)，第VII页。下文引用依据此版。}魏塞(Chr. Hermann Weisse)、卢格(Arnold Ruge)、康士但丁·勒斯勒尔(Constantin Rößler)和布鲁诺·鲍威尔(Bruno Bauer)都有同样的论断；至于问题的实质，即将问题归结到国民经济学，则从黑格尔法哲学出发的那些关于“自然法”和“国家学”的最早努力：贝塞尔(Konrad Moritz Besser)的《自然法体系》（1830年）和艾泽伦(J. Fr. Gottfried Eiselen)的《国家学体系手册》（1828年），都在甘斯的框架下对市民社会进行了描述。埃德曼(Johann Eduard Erdmann)在其《国家讲座》（1851年）中对市民社会地位的理解也差不多。在他看来，市民社会的“本质”在于“明智的计算、理性的考量、工业等”以及“信用”，并且它的这些突出特征首先体现在“地地道道的市民”即“实业家”身上。\footnote{《国家讲座》(Vorlesungen über den Staat, 1851)，第28-29页，及《学术生活与学术研究讲座》(Vorlesungen über akademisches Leben und Studium, 1858)，第97-98页。}但由此埃德曼使市民社会更加靠近中产阶级；这样，黑格尔对“市民社会”一词的使用便与17、18世纪以来欧洲市民等级的解放运动平行。

马克思以及当时的其他人，如里尔(Wilhelm Heinrich Riehl)、洛伦茨·施泰因(Lorenz Stein)、J. C. 布伦奇利(Bluntschli)和瓦格纳(Hermann Wagener)，基本上都是这样理解市民社会的概念的。如马克思在《〈政治经济学批判〉序言》中所言，黑格尔将“物质的生活关系的总和”，按照“18世纪英国人和法国人的先例，概括为‘市民社会’”。\footnote{马克思：《政治经济学批判》，序言，引自1951年柏林版，第12页。（中译文引自《马克思恩格斯选集》第二卷，人民出版社2012年版，第2页。——译者）}提到18世纪英国人和法国人，显然是要强调这个概念是一个特别现代的概念。甚至在马克思之前，就有人（如1850年之前德国最重要的国家学家之一毛伦布雷歇尔[Romeo Maurenbrecher]的做法\footnote{《当代德国国家法原理》(Grundsätze des heutigen deutschen Staatsrechts)，第2版，未修改版（1843年），第22页以下。}）为了在历史上对黑格尔的市民社会进行归类，即由此出发将其与自然法的个体主义观念联系起来。同样，在拉萨尔(Ferdinand Lassalle)看来，霍布斯的“一切人对一切人的战争”(bellum omnium contra omnes)，作为这种自然法的版本之一，不过就是“人们相较于国家，称之为进行自由竞争的市民社会领域”。滕尼斯(Ferdinand Tönnies)也并不是最后一个如是理解的人。他在《共同体与社会》（1887年）中从与现代自然法相区分的概念对立出发，在黑格尔的意义上将“市民社会”与现代的契约、交换和劳动社会联系起来。\footnote{拉萨尔：《巴斯提亚特-舒尔策·德里奇先生》(Herr Bastiat-Schulze Delitzsch, 1864)，引自《演说与著作全集》(Gesammelte Reden und Schriften, ed. Eduard Bernstein)第5卷(1919年)，第69页。滕尼斯：《共同体与社会》，第2版（1912年），第25页。（滕尼斯的著作为《共同体与社会》，正文原文作《社会与共同体》。现改正。——译者）}

在近来的解释家中，罗森茨威格在其著作《黑格尔与国家》（1920年）中追问的不是该词的历史起源，而是其传记上的发生学起源。其结果是指出了黑格尔早期对舒尔策(Sulzer)《一切科学简论》（1759年）以及黑格尔对弗格森《文明社会史论》（1767年）的知识，但这并不能令人完全信服。只有在洛维特这里，黑格尔法哲学的市民社会似乎才第一次得到全面考察，因为他将市民社会与国家的区分（按照他对他援引的《法哲学原理》第260节的具体解释方式）与古代和基督教的二元论、卢梭的《社会契约论》和柏拉图的《理想国》联系起来。与之同时，马尔库塞也能够在其《理性和革命：黑格尔与社会理论的起源》（1941年）中指出内在于黑格尔法哲学之中的社会理论及其对晚近社会学的影响。当然，这必然导致对黑格尔本人及其时代生活于其中的历史维度的论述着墨不多。也正是这种简化处理首先导致了卢卡奇的《青年黑格尔与资本主义社会问题》（1938年写成，1948年出版）一书中解释者与其对象总是不成比例，这个问题唯有成功的个别分析才能解决。当卢卡奇将“市民社会”的表达引入黑格尔的分析（尽管这一表达在他解释的著作中一次也没有出现）时，这一点就最清楚地表达了出来。最后，也能在里特尔的关于《黑格尔与法国大革命》（1957年）的研究中指出同样的问题。此著在现实问题的强有力推动下，强调黑格尔的市民社会第一次从概念上阐述了现代劳动社会。\footnote{参见罗森茨威格，《黑格尔与国家》第2卷(1920年)，第118页；卡尔·洛维特，《从黑格尔到尼采》（1941年），第二部分，第1章，第2节；卢卡奇，《青年黑格尔》，柏林1954年版，特别参见第152页以下、227页以下、280页以下、368页以下以及435页以下；约阿希姆·里特尔，《黑格尔与法国大革命》（1957年），特别是第35页以下、38页以下。此外，参见以下对市民社会的同样重要的解释：魏尔(E. Weil)，《黑格尔与国家》(Hegel et l'Etat, 1950)，第71页以下；罗西(M. Rossi)，《马克思与黑格尔的辩证法》(Marx e la dialettica hegeliana)第1卷《黑格尔与国家》(Hegel e lo Stato, 1960)，第3篇，第543页以下。}

然而，这里出现的问题是，一般说来，这种事后的解释是否抓住了黑格尔的概念的真正历史内容。当我们从后来的社会学角度出发解读黑格尔的“市民社会”时，难道不应该反问一下：它与在康德和沃尔夫那里反复遇到的这个词的以前的用法有何关联？还有：超出康德和沃尔夫以外，这个概念与亚里士多德的政治共同体(koinonia politike)和在西塞罗那里、在中世纪经院哲学和近代自然法当中的市民社会(societas civilis)有何关系吗？因为历史理解的可能性正好在于：将发问背后变动不已的现实——在我们这里是：现代社会现象——一起纳入对过去的理解之中，同时悬置过去自身的视域，并对这些视域进行修正。因此，首先在欧洲政治哲学的概念传统之下对黑格尔的“市民社会”进行打量，就是正当的了。

\section*{1}

在“市民社会”概念史中，一直都有一个值得注意的现象，这就是，这个词差不多同时（1820年左右）不仅以全新的含义出现在黑格尔的《法哲学原理》中，而且也作为政治的核心概念出现在哈勒的《国家学的复兴》中。对于公认的复辟时代理论家哈勒来说，“市民社会”是启蒙哲学的一个革命性概念。启蒙哲学从罗马共和主义的语言中接过了这个词语及其古代含义。这个古代政治概念的人为设立构成了18世纪大革命及其颠覆一切自然—社会关系的真正原因：铲除人与人之间的统治和依附关系、人们服从于国家、消灭世俗和教会的“统治”、消灭自治的团体和同业公会。因此，哈勒在其六大本复辟著作的纲领性前言中宣称，革命时代“所有错误的由来和根源”首先就是“人们将其照搬到完全不同的社会关系当中的罗马市民社会(societas civilis)的不幸观念”。\footnote{参见《国家学的复兴》(Restauration der Staatswissenschaft, 1816-1820)第1卷，前言，第XXVIII页。}哈勒不厌其烦地重复，由于将人们视为一个“市民社会”的自由、平等的成员，因此，“社会联系”也就瓦解了。在哈勒看来，只有自格劳秀斯和霍布斯以来的近代自然法才知道一个与“自然社会”不同的“所谓的市民社会”，他主要是在近代欧洲学术界的拉丁语使用中看见了这种区分的萌芽。\footnote{参见《国家学的复兴》第1卷，第89-90页：“因为拉丁语只有共和主义的言说方式和名称，或者至少在言及国家时，这是最常见的方式和名称，于是这些表达便用到了完全不同的事物和关系上面。同样，正如罗马的市民彼此构成了一个共同体、一个市民团体、一个真正的市民社会：因此，所有其他的人类结合和关系也就必须同样称为市民社会或市民联合体。”}

黑格尔也在《法哲学原理》第182节界定“国家”和“市民社会”时对近代自然法持批判立场。他指责近代自然法没有将这两个概念明确分开：“如果国家仅仅被设想为不同的个人的统一体，亦即仅仅是一个同质的统一体，则其所指就只是市民社会的规定。许多近代国家法学家都不能对国家提出任何别的见解。”针对国家和市民社会的这种混淆和相互混同，黑格尔将市民社会规定为处于国家和家庭之间的“差别”：“市民社会是处于家庭和国家之间的差别[阶段]。尽管其形成要晚于国家。因为它作为差别，以国家为前提，它要存在，前面就必须要有一个独立的国家。”\footnote{《全集》第8卷，第182节，补充，第246页以下。}这里，在市民社会与国家以及家庭与市民社会之间的差别建立起来之后，哈勒所言的“社会联系”也同样瓦解了。但与狂热的复辟理论家哈勒不同，黑格尔肯定地将这一过程规定为国家从社会、社会从国家解放出来的过程，它使国家和社会两者第一次处于它们的真正关系之中。因为他对近代自然法的批判反对显然意味着：当国家与市民社会总是混淆在一起时，国家就不是国家；同样，当市民社会是“政治”社会（即国家）时，它也不是“社会”。

然而，黑格尔关于市民社会及其所谓的与国家混淆的观点仍然停留在哈勒的视域中；两人均将这个概念的起源和历史意义局限于近代自然法，特别是18世纪自然法，而不知道其古老的欧洲传统及其真正的发源地，首先是亚里士多德的《政治学》以及中世纪和近代对亚里士多德的接受。然而，黑格尔自己都不知道，根本而言，他反对的正是这种在近代自然法中依然有效地将国家(polis或civitas)界定为与市民社会(koinonia politike或societas civilis)同一的古典政治哲学传统。因为按照古典政治哲学传统对政治规制的人类世界的理论理解，国家既不自身包含社会，也不预设社会为其存在的前提；相反，国家就是“社会”，然而却是“市民社会”“政治社会”。

“社会”(Gesellschaft)一词古老的原始含义为：结合、联合——希腊语的“koinonia”，拉丁语的“societas”或“communitas”。这样，19世纪以来的、在我们今天看来不言而喻的国家和社会的分裂以及二者的并称，在古代欧洲政治哲学中并不存在。直到18世纪中叶，国家和社会在政治上由同一个概念连在一起，被概括为一个社会整体。这个概念就是古典政治的基本概念“koinonia politike”或“societas civilis”（以前的市民社会）。在亚里士多德《政治学》开头，它将古代城邦的政体表达了出来。这个市民社会从一个自身就是“市民的”社会（也就是政治的社会）的社会实体中接过了其公共政治结构。为此，我们摘引由J. G. 施洛瑟翻译的亚里士多德最早德译版《政治学》第1卷（《政治学》第1卷，第1章，第1节，1251a1—7）译文。尽管译文显然想要在革命时代的经验之后重铸这部政治学的影响力，\footnote{施洛瑟为其翻译写了一个前言，而非导言，它以这个典型句子开始：“值此人人相信自己皆有责任议论和反驳国家形式与革命、市民权利与君主义务之时代，于我而言，将留存至今的几千年前亚里士多德的政治著述译成德语，昭之于世，诚非无益之事。”（亚里士多德：《政治学及经济学残篇》，施洛瑟译自希腊文，附注释和文本分析，第一部分[1798年]，第四页。）}但“社会”一词仍然原封不动在它以前的、与“koinonia”相一致的意义上使用：

\begin{quote}
Es ist offenbar, das ein jeder Staat aus einer Gesellschaft besteht. Eine jede Gesellschaft hat aber, wenn sie sich verbindet, die Absicht, einen gewissen Vorteil zu erreichen. Denn alle Menschen handeln bloß, um das zu erreichen, was ihnen nützlich scheint. Es ist also kein Zweifel, daß die Gesellschaften alle in dieser Absicht zusammentretten, und daß die wichtigste und vortrefflichste, nämlich der Staat, oder die bürgerliche Gesellschaft, auch auf den höchsten und vortrefflichsten Vorteil hinzielt.\footnote{《政治学及经济学残篇》第1页。同样参见施洛瑟之后不久克里斯蒂安·伽尔韦的《政治学》译本（《亚里士多德的政治学》，克里斯蒂安·伽尔韦译，G. 许勒伯恩[Fülleborn]出版、注释并撰文，布雷斯劳，1799年）。此译本同时使用了“市民社会”和“市民联合”。唯有在后来从经济上进行界定的社会概念（尤其是滕尼斯区分社会和共同体）的影响下，晚出的政治学翻译（亚里士多德：《政治学》，E. 罗尔费斯[Rolfes]翻译并撰写导言和注释，1912年；以及O. 纪贡[Gigon]翻译的亚里士多德：《政治学和雅典人的国家》，1955，第1、55页）用共同体(Gemeinschaft)翻译“koinonia”，其后便鲜有例外，而“koinonia politike”也被译为“国家共同体”。在此中立化翻译中，施洛瑟和伽尔韦仍受其约束的古代传统概念的影响力及其典范性已然失效。}（显然，每一个国家都由一个社会组成。然而，当每一社会联合起来时，都是想要达到一定的利益。因为所有人行动，都只是为了达到他们所认为的利益。因此，毫无疑问，所有社会都为此目的而集合在一起，并且最重要、最卓越的社会，即国家或者市民社会，也就指向最高的和最卓越的利益。）
\end{quote}

国家或者市民社会（$\pi\acute{o}\lambda\iota\varsigma\ \kappa\alpha\grave{\iota}\ \kappa o\iota\nu\omega\nu\acute{\iota}\alpha\ \pi o\lambda\iota\tau\iota\kappa\acute{\eta}$；civitas sive societas civilis sive res publica）——这正是古老欧洲政治哲学的经典公式。从亚里士多德到大阿尔伯特、托马斯·阿奎那和梅兰希顿，以及从博丹到霍布斯、斯宾诺莎、洛克和康德，在国家和社会的现代分离之外，这个公式一直有效。大阿尔伯特在他的《亚里士多德评注》中说：“国家、市民沟通或者政治沟通，属于人的自然，亦即那些对人来说乃是自然的事物。”(De numero eorum, quae sunt natura homini, id est, quae sunt naturalia homini, civitas est, et communicatio civilis, sive politica.)\footnote{大阿尔伯特(Albertus Magnus)：《亚里士多德政治学八卷评注》(Commentarii in Octo Libros Politicorum Aristotelis)第1卷，第1章，引自《全集》第8卷(1891年)，第6页。}博丹同样将国家(res publica)规定为市民社会：“由于国家就是市民社会，它无须行会和同业公会即能自立，但它却不能没有家庭。”(Est enim res publica civilis societas, quae sine collegiis et corporibus stare per se ipsa potest, sine familia non potest.)\footnote{《论国家》(De Republica)第3卷，第7章，引自第4版(1601年)，第511-512页。}梅兰希顿和斯宾诺莎只是在形式上偏离了“civitas”的规定，从而将“civitas”和“societas civium”等同起来：“像这样坚实地建筑在法律上和自我保存的力量上的社会就叫作国家，而在这国家法律下保护着的个人就叫作市民。”(Haec autem societas legibus et potestate sese conservandi firmata civitas appellatur, et qui ipsius iure defenduntur cives.)\footnote{《伦理学》，第四部分，命题三十七，附释二。（中译文引自《伦理学》，第200页。译文略有改动。——译者）}梅兰希顿在其《亚里士多德政治学评注》中的说法高度相似：“国家就是一个建立在以彼此利益和最大安全为宗旨的法律之上的市民社会，市民就是那些在这个社会中能够担任公职或者司法职务的人。”(Civitas est societas civium iure constituta, propter mutuam utilitatem, ac maxime propter defensionem. Cives sunt qui in eadem societate magistratus aut judicandi potestatem consequi possunt.)\footnote{梅兰希顿(Melanchthon)：《亚里士多德政治学评注》(Commentarii in Politica Aristotelis, 1530)，引自《梅兰希顿全集》(Ph. Melanchthonis Opera quae supersunt Omnia)第16卷，1850年，布雷特施奈特、宾德塞尔(Bretschneider/Bindseil)编，第435页。}国家和市民社会的传统等同至少在观念上仍然出现在霍布斯的“市民结合”(unio civilis)中，通过让市民服从于国家，市民结合终结了自然状态。对此，他说道：“这样形成的结合叫作国家或者市民社会。”(Unio autem sic facta appellatur civitas sive societas civilis.)\footnote{霍布斯：《论公民》，1642年，第5章，第9节。}最后，在这一脉络中，我们还可以提到洛克的《政府论》（1689年）和康德的《道德形而上学》（1797年）的第一部分。洛克此书第7章的标题是“论政治社会或市民社会”，而康德也在第45节将国家界定为“civitas”，并在第46节中将其视同于被解释为“societas civilis”的国民社会(staatsbürgerliche Gesellschaft)：“这样一个社会(societas civilis)（亦即一个国家）的成员，为了立法的目的集合起来，就叫国民(cives)。”康德意义上的“国民”(Staatsbürger)——这个词被他的同时代人当作一个新词，一开始还遭到拒绝——事实上就是旧的市民社会的具有法律能力和政治权利的市民，即城邦的市民(politai der polis; cives der civitas)。他们具有特别突出的国民品质，“了解市民事务和审慎”(“rerum civilium cognitio et prudentia”)（西塞罗：《演说家》第1卷，14，60），参与公共生活、立法和国家管理。因为直到18世纪，只有生活在市民状态或政治状态(status civilis sive politicus)中的那些人才是“市民”，市民社会建立在这种状态之上。它不是“社会”观察方式的主题，而仅仅是广义上的政治学的主题。沃尔夫说道：“对生活在国家或者市民状态下的人进行思考的哲学部门，叫政治学。因此，政治学就是指导市民社会或者国家中的自由行动的科学。”(Ea philosophiae pars, in qua homo consideratur tamquam vivens in Rep. seu statu civili, Politica vocatur. Est itaque Politica scientia dirigendi actiones liberas in societate civili seu Rep.)\footnote{克里斯蒂安·沃尔夫(Christian Wolff)：《理性哲学或逻辑学》(Philosophia rationalis sive Logica, 1740)，预备性讨论，第3章，第65节。}

因此，我们可以说，“市民社会”在古老的欧洲意义上是一个传统的政治概念，是政治世界的基本核心范畴，其中“国家”和“社会”尚未发生分离，反倒构成了市民—政治社会自身同质的统治结构。这种统治结构建立在家庭奴隶劳动、奴隶制度或者农奴制度和工资制度的“经济”领域之上，并与之形成鲜明对照。因为在政治学的这个古典传统中，并非共同体的所有居民都拥有市民身份；所有类型的非自由人：那些居于公共政治的市民领域之下、在家庭的私人领域中为了基本生存需要而劳动的人，以及另一方面，同样从事“经济”活动的、束缚于家庭作坊的手工艺人以及妇女，都不属于市民社会或国家(societas civilis sive res publica)，因为他们是“Oikos”（即“家庭社会”）的一部分，没有赋予他们市民身份的政治地位。亚里士多德政治学的原理，同时构成了其经济学的前提：“奴隶没有城邦”（《政治学》第3卷，第9章，1280a32）。正如康德在其时代（18世纪末）所言，这一原理也依然适用于佣人、日工以及作为佣工的手工艺人。他们生活在真正意义上的市民社会之外，市民社会参照他们得到界定，并与他们形成对照。除了从属于政治以外，这一点构成了旧市民社会的第二个基本结构要素：往下与经济、与家庭分离——市民社会(societas civilis)与家庭社会(societas domestica)的古典政治对立。在托马修斯、沃尔夫和康德那里，都还能找到这种对立。托马修斯说：“然而人的社会本身要么是市民社会，要么是家庭社会。家庭社会构成市民社会的基础，因为这里市民社会的意思不外乎许多家庭社会及其人口的联合，而处于一个普遍的统治之下。”\footnote{托马修斯(Chr. Thomasius)：《政治明智短论》(Kurtzer Entwurf der politischen Klugheit, 1725)，第204-205页。市民社会与家庭社会之间的这种对比也同样适用于康德，因此，尽管康德提出了国家革命的、着眼于市民自由的新原则，但仍然保持了市民社会与国家的古老同一，因为他仍然知道旧欧洲的“家庭法”（《道德形而上学》，第一部分，1797年，第24节以下。引自科学院版，第6卷，第276页以下），以及与之相联的“物的方式的人格权”（同书，第22节以下，第276页以下），由此导致了他的全部困难，即要规定一个“构成他自己的主人(sui juris)”的“人的等级”（见“论俗语：这在理论上可能是正确的，但在实践上是行不通的”[1793年]，第二部分，引自《全集》[科学院版]，第8卷，第295页）。}

\section*{2}

历史政治生活的细节复杂多变，从亚里士多德到康德的伟大政治著作对它们的书写也各不相同。如果我们不考虑这些，那么就得承认，就国家与社会的现代关系而言，直到黑格尔以前，古老的欧洲学术传统中，没有它的“概念”，而根据传统的“市民社会”概念，这一点也绝非偶然。这个概念的连续性，它传承的漫长历史及其对于政治学理论所具有的典范意义（在黑格尔这里，这种典范意义似已不再存在），也并不能阻止黑格尔进一步使用这一概念，与此同时又与这个传统以及古老的政治世界决裂。黑格尔完成了这个传统和这个古老的政治世界，表示与之决裂，并且提出了一个新的市民社会概念。只有在1800年左右这个伟大的欧洲传统发生断裂的大背景下才能够说明，为什么黑格尔这个欧洲形而上学中的历史思想的代表，在历史性地反对市民社会及其与国家的“混淆”时，仍然停留在哈勒的视域中。就事物本身而言，我们也可以反过来说，要对这个概念做出新的界定，就必须跟古老的欧洲政治传统一刀两断。18世纪末的革命开始了这一断裂，如在哈勒处所见，其中也存在着激烈反对这种趋势的复辟思想。唯有在这种语境下阅读《法哲学原理》中黑格尔市民社会概念的引入，并像这里一样反对从马克思和洛伦茨·施泰因出发对它所做的流行的历史性解释。\footnote{这一方面最典型的是福格尔(Paul Vogel)：《黑格尔的社会概念以及L. 施泰因、马克思、恩格斯和拉萨尔对它的历史续造》(Hegels Gesellschaftsbegriff und seine geschichtliche Fortbildung durch L. Stein, Marx, Engels und Lasselle, 1925)。参见历史学家布伦纳(Otto Brunner)对现代社会概念问题的（哲学上有关的）与之相对的阐述，载《社会史的新道路》(Neue Wege der Sozialgeschichte, 1956，尤其是第7页以下、30页以下和207页)；康策(Werner Conze)：《德国早期革命时代的国家和社会》(Staat und Gesellschaft in der frührevolutionären Epoche Deutschland)，载《历史杂志》(Historische Zeitschrift)第186卷(1958年)，第1页以下，此文扩展后收入文集《德国三月革命前的国家与社会》(Staat und Gesellschaft im deutschen Vormärz, 1962)，第207页以下，我从此文中获益良多。}换言之，首先从政治传统和前革命世界出发进行解读，才能衡量它对于19世纪所具有的真正意义和历史现实性。

从亚里士多德直到康德的政治形而上学大传统将国家称为市民社会，因为在市民的、康德依然用拉丁语对其进行解释的“cives”（市民）的法律能力中，社会生活本身就已经是政治的，而且这个如此理解人类世界的政治状态(status politicus)，同时也将处于统治—家庭的或者等级的分层中真正“经济的”和“社会的”要素涵括于自身当中；反之，黑格尔则将政治领域、国家同一——现在已经成为“市民的”(bürgerlich)——“社会”领域分开。由此“市民的”这一表达便对立于其原始含义，获得了一种首要的“社会”内容，从而不再像18世纪那样，在与“政治的”相同的意义上使用。当国家在政治上成为绝对，并由此出发保证社会拥有它自己的重心，将其作为“市民社会”释放出来的时候，这种表达便仅仅指称在这种国家中私人化了的市民的“社会”地位。当人们就像自从马克思、拉萨尔、滕尼斯以来所做的那样，在近代自然法中追溯黑格尔的[社会]概念的起源，就会将这一历史过程变成自然状态(status naturalis)和市民状态(status civilis)的自然法模式，而没有正确对待近代自然法本身以及在其中仍然发挥作用的传统。因为譬如说霍布斯的市民社会(societas civilis)（《论公民》，第1章第2段；第5章第9段）也并不是那种“bellum omnium contra omnium”（一切人对一切人的战争）、彼此之间互相契约的交换社会（如亚当·斯密和重农学派最先设想的）；相反，它依然终究是一个政治秩序概念、自然状态(status naturalis)的对立物。在德国，直到施洛策(August Ludwig Schlözer)为止\footnote{参见施勒策(August Ludwig Schlözer)：《国家法总论》(Allgemeines Stats-Recht, 1794)，其中市民社会成员被称为“宗子，因而是自由的家长”（第一部分：《元政治学》[Metapolitik]，第17节，第64、65页）。施勒策用传统公式表达了它的结构：“Societas civilis或者civitas”，以区别于作为“Societas civilis cum imperio”的“国家”（同书，第44节注释）。他在这种仅仅暗示了黑格尔的“差异”的对立中，把握到了“国家与市民社会”之间的区别。}，亚里士多德的独立家长的观念都一直保存在自然状态(status naturalis)中，这种家长在家中自由地发号施令，因此属于市民社会的成员。在施勒策这里，市民社会第一次与“国家”形成鲜明对照，这一事实乃是独立于自然法模式、基于前述国家与社会之间的历史“差异”而发生的。随着政治集中到国家（行政、宪法、军事）、“市民事务”转移到社会，便同时产生了这种历史差异。正是黑格尔命名的这种差异规定了1821年《法哲学原理》“市民社会”部分的内部结构和历史内容。这里，从其政治—法的意义中解放出来的市民概念便与同时得到解放的社会结合起来了。当旧的市民社会(societas civilis)经典表达的政治实体消解为社会功能时，“市民”(Bürger)和“社会”两者，通过18世纪末对欧洲传统的突破，随着工业革命和政治革命而获得了这种社会功能。至此，市民才作为布尔乔亚，成为政治哲学的根本问题，这是与现代社会的形成同时发生的。由于现代社会自身基本上接过了“经济”的功能，便一天天消解了旧的家族实体。以前，家庭是旧的市民社会的社会细胞，现在，则是市民社会的变化形态构成了现代国家的社会基础。

在政治概念的内在转变中，这里简单勾勒的这个历史过程更加具体地显露出来。1800年以后，“市民社会”的结构似乎在一个汲取了旧的意义成分的新的中心上面建立起来。因此，只要我们注意到自然法理论在人作为人与人作为公民之间所作的区分（北美革命和法国革命的人权与公民权将这种区分法典化了），那么黑格尔的概念与18世纪的概念之间的区别及其内在地属于19世纪也就一目了然了。按照自然法思想，作为人(Mensch)，他是“societas generis humani”（人类社会）的成员，同时是类存在(Gattungswesen)和个体，服从伦理学法则（这种法则在其普遍性中是没有规定的）。作为公民，他属于市民社会、国家及其法律，服从政治学所要求的规则。\footnote{对此，参见克里斯蒂安·沃尔夫：《理性哲学或逻辑学》（1740年），预备性讨论，第3章，第64-66节；康德：《道德形而上学》，第一部分（1797年），法权论导论；卢梭：《爱弥儿》（1762年）第3卷，见《全集》第3卷（1823年），第371页，等等。}相反，在《法哲学原理》的“市民社会”中，自然法上作为人的人，也就是类的代表，消融于他的自然需要、消融于黑格尔在第190节解释中嘲讽般地谈到的“观念的具体物、人们所谓的人”及其本身归于“需要的立场”之下的东西。\footnote{《全集》第8卷（1833年），第256页。参见接下来的限制（我们现在必须参照马克思对需要体系的普遍化以及将人固定在生产和消费之间的运动上，阅读这种限制）：“因此，只是在这里，并且真正说来也只在这里，才在这种意义上谈到人。”（第256页）}这个具体物，第190节基于人和动物的自然区分将其规定为抽象—普遍的需要者，因而现在就是市民社会的对象，即“市民”——黑格尔用法语在括号中标明为“bourgeois”。作为单纯的（即自然的）人，人就是需要者，作为需要者，他就是私人，即作为“bourgeois”的市民。人和市民不再像在18世纪中那样彼此对立；相反，在现代市民社会，市民将人包含于自身。

在历史概念的这个改变了的结构中，还隐藏了一个进一步的联系：欧洲革命与古老的古今之争之间的联系。这里，这个联系只是通过暗示得到了表达。18世纪末的一个真正历史性标志就是，在政治学中，正是通过接受古人的政治概念，开始意识到与古代、与希腊人和罗马人的古老的市民自由之间的距离。因为在欧洲政治传统的历史中，这种接受第一次无法“应用”\footnote{这个概念在这里的含义，在伽达默尔1960年出版的《真理与方法》一书中所指出的解释学的传统和本质环节的意义上使用。}于时代的历史实际：现代领土国家、“臣民”、城市居民、现代社会等。这种变化明显出现在18世纪的市民概念的问题上，也同样出现在我们即将看到的新“社会”给政治思考所带来的问题上。骤然间出现的旧概念无法应用的情况，其最典型的表征可能是卢梭的建议：将诸如“公民”(citoyen)或是“祖国”(patrie)之类的词语完全从语言中抹去。\footnote{参见卢梭，《爱弥儿》（1762年）第1卷（引自《全集》第3卷，1823年，第15页）。（中译文参见〔法〕卢梭，《爱弥儿》上卷，李平讴译，商务印书馆1978年版，第11页。——译者）}那时，德国的思想没有这样激进，人们也不知道“bourgeois”（市民）和“citoyen”（公民）之间的区别，而是满足于这样的定义：市民概念在一种“普遍的意义”和一种“特殊的意义”上得到界定。市民作为“国家的成员(civis)”，属于前者，由此，与传统模式一致，“君主自身”出现为“民族的第一市民”。\footnote{对此，参见当时影响很大的哥廷根法学家谢德曼特尔(Scheidemantel)，《德国国家法和采邑法》(Teutsches Staats-und Lehnrecht, 1782)，第一部分，第2节。同样，阿亨瓦尔(Achenwall)也说，邦君并不完全在市民社会以外，而是“要把自己当作国家最杰出的市民(civis eminens)”，见《欧洲帝国的国家宪法》(Staatsverfassung der europaischen Reiche, 1752)，第2页。两个人的真正问题是市民和臣民的关系，因此就是市民的政治概念，因为“我们发现有许多臣民不是市民，比如说外邦人或者奴隶；其他一些人既是臣民也是市民；相反，也有不是臣民的市民，如君主。”（《德国国家法和采邑法》，第一部分，第439页）。}人们所理解的“特殊的意义”下的市民，则是城市成员、聚居区居民(Ortsbürger)，因此就是有别于贵族和农民的市民等级。现在，决定性的是，由于现代政治国家几乎铲除了市民(civis)的社会基础、市民社会(societas civilis)当中被统治者的自治，因此也就将市民概念的“普遍意义”减损至最低。黑格尔在1800年左右就已经罕见地理解了这一过程；因为只有从他开始，才直接出现了公民(citoyen)与市民(bourgeois)、私人与扩展到所有“臣民”的国民(Staatsbürger)的并列。现在，关心“普遍物”的政府、国家与“以个人为目的”的另一“端”即“市民”之间的现代差异同“古代伦理”及其政治实体形成了鲜明对照：“同一个人既关心自己和家庭、劳动、订立契约等，同时他也为普遍物劳动，以普遍物为目的。按前一方面，他叫市民，按后一方面，他叫公民。”\footnote{黑格尔：《耶拿讲座》（1805-1806年），J. 荷夫迈斯特编，《耶拿实在哲学》第2卷（1931年），第249页。黑格尔还加了边注，将这些词语转译到特殊的德国语境：“小市民和帝国国民，两者形式上都是小市民。”}青年黑格尔已经在希腊城邦共和国的衰落中瞥见这一过程，并且想要在罗马法中确立市民的法律地位。毫无疑问，这里，在晚期古代的那些现象中，他尖锐地道出了他的时代问题。

在晚期柏林讲座中，黑格尔在分析亚里士多德《政治学》的最后，以史论政，比较了现代国家和唯一与之相符的现代社会现象：工厂。他在这里谈到，“我们近代国家的抽象权利把个人孤立起来，允许他如此”，继而建立起一种必然的联系，以至于“真正说来，没有人具有整体的意识，或者为了整体而活动；他为整体而工作，但是不知道他是怎样为它工作的，他只关心他个人的保护”。这里，黑格尔指出了近在眼前的现代工厂形象。在现代工厂，就像在国家和社会中一样，整体活动的生命分散开来，分成不同的部分：“这是一种分工的活动，在这种活动中，每个人只占一份；正如在一个工厂里面，没有人自己单独制造一件完整的产品，每个人都只制造产品的一部分，不懂得制造其他各部分的技能，只有少数几个人才把各部分装配成一件产品。”只有古代产生的那些自由人民才有整体的意识，并且为了整体而活动，现代人作为个人本身就是不自由的——“市民的自由就是不需要普遍的事物，就是孤立的原则。但是，市民的自由（我们没有不同的两个词语来表达bourgeois[市民]和citoyen[公民]）是一个必然的环节，那是古代国家所不知道的。”\footnote{《哲学史讲演录》第2卷，K. L. 米希勒编辑，引自《全集》第14卷（1833年），第400页。（中译文参见《哲学史讲演录》第二卷，第364-365页。——译者）}

\section*{3}

然而，黑格尔直到晚年才完全洞见这个环节的历史必然性及其积极价值；因为毫无疑问，他只有在1820年左右才通过《法哲学原理》的市民社会部分“把握到”欧洲社会和政治制度内部的这种根本变化，并且这种“把握”恰恰来自于他对这种变化所做的系统的、富有成果的论述，而不像他在伯尔尼、法兰克福和耶拿开始经济学和政治学研究之后，仅仅对其进行断言而已。只要浏览一下他耶拿时期的文章、摘录和讲座，就会立即发现，根本而言，他从英国人和法国人的著作研究中所得到的事实，以及他对现代社会制度的理论洞见，同他对古代文化遗产的创造性接受漫无关联地混杂在一起。由此，如前述，市民(bourgeois)与罗马帝国时期并置，同时与希腊“政治生活”(politeuein)的“绝对伦理”相对立，从斯图亚特、斯密和李嘉图那里摘录的社会、需要、劳动和占有的经济领域获得了“相对伦理”的名称，并且作为古代悲剧和喜剧的背景出现，而对一般劳动制度的分析要么直接在“人民”和“政府”的标题之下进行，要么彼此独立，出现于人与自然的“实践”过程的框架之下。\footnote{对此，参见《批判的哲学杂志》第2卷，第2期（1802年），第79页。《政治学与法哲学文集》，G. 拉松编，第418页以下、464页以下，首先参见1803-1804年《耶拿讲座》（《耶拿实在哲学》第1卷，J. 荷夫迈斯特编，1932年，第220页及以下、236页及以下）和1805-1806年《耶拿讲座》（《耶拿实在哲学》第2卷，1931年，第197页及以下、231页及以下）。}除了他倾慕不已的古人鲜活的文化传统以外，黑格尔也直接指出了劳动和分工、机器、货币和商品、贫与富这些现代人的铁的事实，但未能深入它们的概念之中，或者至少未能将它们系统地结合起来。在耶拿时期，黑格尔似乎拒绝将这个凌乱世界的丰富内容结合起来，建立起他后来称之为人类世界的历史领域的“客观精神”的类别。另一方面，当他首次在1817年《哲学全书》初稿中试图结合古代和现代而建立这种内容方面的类别时，概念构造的古典化风格令他总是无法对现代状况做出到位的阐述，特别是在与该著1827年第2版进行对比时，这一点就非常明显。唯一论及这一主题的海德堡《哲学全书》第433节只谈到“普遍作品”及其特殊化为“诸等级”。尽管家庭表现为“个别性等级”，但是其作品为“特殊定在的需要，而其切近的目的是特殊的主体性，但是，要达到这一目的则以其他一切人的劳动为前提，同时也会影响到这一劳动”的领域，还不是后来的“市民社会”，而是“等级”，更确切地说，是（如其模糊隐约所言）“特殊性的等级”。\footnote{《哲学全书》（1817年），第433节。（中译文引自《哲学科学全书纲要》，第309页。——译者）对此，可以参考黑格尔于《法哲学原理》出版之后在《哲学全书》第2版（1827年）对这一段所做的重要改动（第513节以下）。这里，同在其他地方一样，由于“伦理”这一术语与近代家庭、市民社会和国家的发展联系起来而历史化了（第552节）。而在耶拿时期末，在《精神现象学》（1807年）中，对古代和现代的最早的历史图式化就已经出现，因此，罗马法状态便位于阿提卡共同体之后（在“真实的精神，伦理”部分），以及国家权力和财富的二元论表现为17、18世纪的特殊现象（在“自我异化了的精神，教化”部分）。}因此，对于社会领域，我们可以说，黑格尔正是通过近代经济与社会的概念结构与希腊城邦生活的概念结构之间的巨大差距，才达到了他对他自己历史地位的认识。正是由于旧政治学领域的古代传统概念无法应用于革命世纪的社会状况，导致黑格尔1820年左右构造出“市民社会”概念，作为国家和家庭之间的差异领域。因为在1820年《法哲学原理》面世以前，他既没有用过这个词语，也没有用过它的概念上的根本含义。因此，罗森茨威格提示的他在学生时代对舒尔策的《一切科学简论》摘录和弗格森的《文明社会史》是不正确的，因为在这两种情况下，讨论的都是旧的传统概念（尽管已非其本来面目）。在罗森茨威格明确提到的地方，舒尔策考察的是旧的意义上的市民社会，用他的话来说，就是“市民国家”的制度。因此，对舒尔策来说，从“市民社会的普遍概念”只是产生了“在上者的权力和司法权、臣服、奖惩之类的特殊概念”。\footnote{舒尔策：《一切科学简论》（1759年），第189-190页。}弗格森同样在与“政治社会”相同的意义上使用这个概念（第三部分，第6节），但在这个标题下，在广泛阐述（首先是斯巴达、雅典和罗马的）政治状况之余，也对“艺术和科学”进行了阐述。

如果这就解释了黑格尔的[市民社会]概念的外在起源的话（此诚为可疑），也依然有问题未获解决：黑格尔为何只在学生时代的摘录和早期神学政治片段中使用这个概念，\footnote{参见《早期神学著作集》（1907年），第41页以下、44、191页。在这些地方“市民社会”(bürgerliche Gesellschaft)依18世纪末的用法，与“国家”同义。就此而言，霍切瓦尔(Rolf K. Hočevar)在《青年黑格尔论等级制和代议制》(Stände und Repräsentation bei jungen Hegel, München 1968；即《慕尼黑政治学研究》第8卷)第9页第23条注释以及第201页第79条注释中对我的主张所提的反驳是站不住脚的。《早期神学著作集》第191页的这个句子特别明显地表明了这个词语的传统用法：“在市民社会的本性中……个人的权利变成了国家的权利。”这里，黑格尔的原型是门德尔松(Moses Mendelssohn)：《耶路撒冷》(Jerusalem, 1783)，引自莱比锡1869年版，第132页以下、146页。}其后却从未在其著作的其他地方用过它？在他的著作中，情况是这样的：即黑格尔与大约自1780年以来在克劳斯(Christian Jakob Kraus)、根茨(Friedrich Gentz)、洪堡(Wilhelm von Humboldt)、胡弗兰德(Hufeland)等人处出现的“societas civilis sive civitas”（市民社会或国家）的古老传统表达的消退甚至消解相一致，用“国家社会”(Staatsgesellschaft)来表述之。\footnote{对此，参见我的论文“市民社会”载《近代政治-社会概念词典》(Lexikon der politisch-sozialen Begriffe der Neuzeit)第2卷（1970年），O. 布伦纳、W. 康策、R. 科泽勒克(Koselleck)出版。}这个措辞非常清楚地表明，他不再使用旧的市民社会的概念，但是也还没有使用他自己的新概念。按照《纽伦堡哲学入门》，这个所谓的国家社会分为家庭和国家，前者是“自然社会”，后者是“人们通过法律关系形成的社会”。两者之间还没有后来所言的“市民社会”，因为“家庭的自然社会扩展为普遍的国家社会，它既是一个通过自然建立的联合，也是一个通过自由意志进入的联合”。\footnote{《纽伦堡哲学入门》，引自《全集》第18卷（1840年），卡尔·罗森克兰茨编，第47页以下、199页。[本段引文原文现收入《黑格尔全集》（历史考订版）第10卷第1分册，第359页。中译文参见《黑格尔全集》第10卷，第292—293页。——译者]}

然而，对于黑格尔的思想世界来说，这些处于《耶拿讲座》和海德堡《哲学全书》之间的纽伦堡中学预备课程，意义非常有限。由于在这些课程中“市民社会”了无踪影，因此它们就只是预兆性的。只有到了1817年以后，也就是抛弃了古代化的政治概念构造以后，黑格尔在耶拿时期获得的实际认识——在古代的政治遗产以外，随着近代“社会”和“市民”(bourgeois)，出现了一种与古典世界对立的不同现象——才成为对现代国家制度内部结构原则性的历史洞见。1820年，黑格尔将他1805-1806年在“个别劳动者”领域所见到的以及（这里，如果愿意的话，也总归是“古典的”）仅仅局限于这一领域的东西——“他（劳动的个人——笔者）在普遍物中有其无意识的实存；社会是他的自然，他依赖于社会的基本的、盲目的运动。这种运动从精神上和肉体上维系他，或扬弃他”\footnote{《耶拿实在哲学》第2卷（1931年），J. 荷夫迈斯特编，第231页。}——在概念上表达为“市民社会”以及对于市民社会来说同样典型的基本必然性。市民社会出现了市民的“市民自由”、他的部分依赖于“体系”的定在，在历史哲学讲座中，黑格尔谈到了这种定在。这样看来，黑格尔将“市民的”和“社会”组合成政治哲学的一个基本概念绝非偶然。这个概念，尽管外在地看，与亚里士多德的“政治共同体”(koinonia politike)、博丹、梅兰希顿或者沃尔夫的“市民社会”(societas civilis)和康德的“市民社会”的传统相一致，但其产生却恰恰是以与这个传统的完全断裂为前提的。就此而言，我们可以完全正当地说，一般而言，在黑格尔以前，甚至在他1820年以前的著作中，都不存在现代意义上的市民社会概念，只有通过这个概念，他才可能以一种影响重大却又必然的、富有成果的方式解决前面提到的“应用”问题。由此，在政治哲学的结构内，自博丹为主权概念奠基和卢梭提出公意以来，黑格尔做出了最为重要的改变。在他之后不久，这一改变即成为现代人的真正问题。

\section*{4}

黑格尔用“市民社会”对时代意识所造成的改变，绝不亚于现代革命之后果：通过政治集中于君主国家（更确切地说，革命国家），产生了一个去政治化了的社会，并且其重心转移到了经济上面。在同一时期，随着工业革命，这个社会在“国家经济学”（更确切地说，“国民经济学”）中经历了这种转变。唯有在这一过程中，欧洲社会的“政治”制度和“市民”制度始告分离。在此前古代政治的古典世界，“政治”制度和“市民”制度的意义相同：如托马斯·阿奎那所言的“communitas civili sive politica”（市民共同体或政治共同体），洛克所言的市民社会或政治社会。黑格尔自己也完全知道市民事务与政治事务的这种古老同一，在《法哲学原理》的一个地方还谈到了这两个词语在原始意义上的统一。他在第303节言及市民等级的、非政治的社会现代观念及其与旧的政治等级差别时，说道：“尽管在这些所谓的理论看来，一般的市民社会等级与政治意义上的等级相去甚远，然而语言却仍然保持了先前本来就有的这种统一。”\footnote{参见《全集》第8卷(1833年)，第303节，第398页。马克思自然会在其对黑格尔《法哲学原理》的批判性评论中注意到这句话，因此他立即写道：“因此，我们可以推论，现在这种统一已不复存在了。”（《黑格尔国家法批判》，1843年，引自《马克思恩格斯全集》，第一部分，第1卷第1册，第487页。）}在这里，黑格尔反对国家与社会对立的自由主义革命观念，但这并不能改变这一事实，即他在概念上比反对国家的自由主义运动还要尖锐，认识到了现代革命与“市民社会”的联系，并将其表达出来。因为黑格尔开始阐述市民社会的第一段，即《法哲学原理》第182节，便从孤立的个人、从19世纪早期在政治制度上得到了解放的社会（因而成为了前述意义上的“市民社会”）的私人这种事实出发，明确表达了当时出现的旧的市民社会传统的瓦解：“具体个人作为特殊的人自身就是目的。他作为各种需要的整体以及自然必然性和任意的混合体，构成市民社会的一个原则，——但这个特殊的个人本质上是与其他的这样的特殊性相关的，从而每个人都通过他人的中介，并且同时也完全只是作为通过另一个原则（即普遍性的形式）的中介，使自己有效并得到满足。”\footnote{同上书，第246页。下面的解释，参见我的论文“黑格尔《法哲学原理》中的传统与革命”，载《哲学研究杂志》(Zeitschrift f. philosophische Forschung)第16卷第2期(1962年)，特别是第218页以下，见本书第170页以下。}通过这一段话，之前已经解放了的西欧社会的功利原则及其出现在曼德维尔和斯密那里的经济模式、狄德罗和爱尔维修的个人利益、边沁和富兰克林的自利便以德国哲学的形式出现了。这里，结果也完全一样：在个人与个人之间，产生了一个独立的、基于利益的关系网络。正如黑格尔在下节所言，市民社会的两个原则“在其实现中的利己的目的”与它的活动条件，即利己的“普遍性”，建立起“一个在所有方面相互依赖的体系”，并且是这样的：“单个人的生活和福利及其权利的定在，都同所有人的生活、福利和权利交织在一起，建立在所有人的生活、福利和权利的基础上，并且唯有在这种联系中，单个人的生活和福利及其权利的定在才是现实的和有保障的。”\footnote{《法哲学原理》，《全集》第8卷(1833年)，第247页。在《哲学全书》第2版和第3版中，黑格尔简单地将市民社会称为“原子论体系”（参见《全集》第7卷第2册，L. 布曼[Boumann]编，1845年，第395页）。}要能够由此出发理解社会之为“市民社会”的描述，我们阅读第183节时，就必须对照第187节。第187节的第一句话是：“个体，作为这种国家的市民，就是以其自身的利益为目的的私人。”然而，这些私人性市民个人的单个利益所追求的目的建立在所有利益目的的内在联系之上，因此要达到单个利益的目的，这些个人就必须按照那种“普遍的方式”规定他们的意向和行为，从而使利益的联系、“社会”直接出现在他们的对面。同时，他们亦如黑格尔所言：“使自己成为这种社会联系之链上的一个环节。”\footnote{《全集》第8卷(1833年)，第251页。参见对这种联系作为一种“社会联系”的描述，它“是所有人都从其中得到他们的满足的普遍财富”。见柏林时期《哲学全书》，第524节（《全集》第7卷第2册，第395页）。}

正是这个普遍依赖的体系及其联系之链，发展出了有别于国家的、从需要和劳动的土壤中生长起来的并通过其活动不断产生的“私人”（第187节）之间的关系网络——“社会”（在该词的现代意义上）。黑格尔在其时代（1815年到1830年之间）的社会架构下（此时，社会尽管依然与等级相连，但已脱离了政治国家），详细阐述了社会的内部划分：经济上的需要的体系（第189-208节）、私法上的司法（第209-229节）以及通过警察和同业公会（第230-256节）实现的国家中的政治伦理整合。这个社会表现为“市民的”，因为它根本上已经逐渐演变为由私法调整的“作为私人的市民”的利益体系；它依然是政治的，实际上只是就它通过等级的和同业公会的团结，在与“警察”的联系中，自我运行在一个相对稳定的秩序体系当中而言的。

因此，我们必须认为黑格尔对市民社会的这种论述（这一论述并非如通常理解所认为的那样特别着眼于普鲁士，而是一般地着眼于1815年后从革命“倒退”的时代\footnote{甘斯在其关于《自然法和普遍法史》(Naturrecht und Universalrechtsgeschichte)的柏林讲座(1832-1833年冬季学期)导言中使用了这一表述。}）是整体性的，也是历史性的。可以说，它是一次伟大的尝试，既不将“社会”建构为基于劳动和享受的需要体系，也不将其建构为一种旧的欧洲社会安于旧的政治状况、安于统治和依附的自然因素，而是对革命以来欧洲人类世界得到解放的“社会”存在与其“政治的”、法的和伦理的秩序进行调解，并在调解中将它们提高到它的新概念上。因此，黑格尔的“市民社会”包含了后来的19世纪“社会”、社会学和国民经济学的对象、如他自己所言的“原子论的体系”，但这只是它的最底层；需要的体系尽管开始了社会的运动，但是社会运动要由“司法”进行规范和组织。因为，诚如黑格尔在第229节所言，在司法中，“市民社会（在市民社会中，理念迷失在特殊性之中，并且陷入内外分裂）回复到它的概念，即自在存在的普遍物与主观特殊性的统一”。\footnote{参见《全集》第8卷，第293页。}如果社会没有在法、伦理和政治上得到规治与团结，就不是“市民”社会。因此，黑格尔此处就司法而言的“概念”，同时也就是市民社会本身的概念。在黑格尔这里，这个概念在新的经济和私法成分（对社会来说，这些成分是构成性的）之上，依然保存了旧的伦理和政治的结构，尽管缩减为“警察”并局限于“同业公会”。因为，另一方面，社会的实体必然性，即警察和同业公会两者对它的“调整工作”（第236节）和“伦理化”（第253节，补充），便在同业公会中将古典世界的实质、“伦理的力量”（第145节）制度化了，并在警察中将“政治制度”（第269节）理性化了。但这也同时表明了黑格尔思想的历史穿透性，即他正好将警察和同业公会这两种秩序要素纳入他对市民社会的阐述之中。

对我们来说，将这两个概念相提并论不仅不恰当，而且更是不可理解：“同业公会”具有古代意味，“警察”则极具实际意味。而如《法哲学原理》所示，黑格尔通过它们描述的事情，并不能立马表明二者之间的内在亲和性及其与市民社会的变化形态之间的历史联系。要指出这种——对黑格尔的特殊概念来说恰好是构建性的——联系，我们就必须透过他的描述，回到直到18世纪为止的强大的古代政治学传统，追问其中的“警察”和“同业公会”究为何物，以及黑格尔本人还在多大程度上意识到它们的结构。首先要说明的是，《法哲学原理》意义上的警察概念跟我们的警察概念关系不大。这个概念最初随着17-18世纪近代国家的制度化和官僚化而得以接受，包含了国家对内部分化开来的社会进行的一般管理。就此而言，它是旧的“市民社会或国家”分化为“国家”和“社会”领域的概念上的对应物。更准确地说，它在管理的手段中标明了去政治化了的社会与政治国家之间的中介。“警察”本身就是变成了国家行政的旧政治，它将市民社会的政治制度及其统治技术表达了出来。\footnote{18世纪下半叶是“警察学”的伟大时期。1830年，青年莫尔(Robert von Mohl)仍以大部头的警察学论述开始其国家法研究。当时的著作有冯·尤斯蒂(J. H. G. von Justi)：《警察学原理》(Grundsätze der Polizeywissenschaft, 1756-1759)、冯·索南费尔斯(J. von Sonnenfels)：《警察原理》(Grundsätze der Polizey, 1765)、冯·普法伊费尔(J. F. von Pfeiffer)：《自然的、产生于社会的终极目的的普遍警察》(Natürliche, aus dem Endzweck der Gesellschaft entstehende allgemeine Polizey, 1779)、勒西希(K. G. Rössig)：《警察学教程》(Lehrbuch der Polizeywissenschaft, 1786)、荣格(J. H. Jung)：《国家警察学教程》(Lehrbuch der Staatspolizeywissenschaft, 1788)、洛茨(J. F. E. Lotz)：《论警察的概念》(Über den Begriff der Polizey, 1807)，等等。参见迈尔(H. Maier)近著《旧德国国家学和行政学（警察学）》(Die ältere deutsche Staats- und Verwaltungslehre [Polizeiwissenschaft], Neuwied/Berlin, 1966)。}当青年黑格尔在1805-1806年耶拿讲座中将现代“警察”同政治的起源、古典希腊哲学的“Politeia”（政制）联系起来时，还清楚地意识到这种联系：“这里，警察来自‘Politeia’（政制）、公共生活和统治、整体自身的行动，然而现在已经降格为整体提供各类公共安全、保护企业反对欺诈的行动。”\footnote{参见1805-1806年《耶拿讲座》，引自《耶拿实在哲学》第2卷，J. 荷夫迈斯特编(1931年)，第259页。}警察就是国家和社会之间的差异得以出现且不断得到调解的形式：它似乎是在市民社会的独立形态中唯一可能的“政治”结构，即行政管理，如《法哲学原理》第235节所言的“公共权力的监督和照管”。\footnote{《全集》第8卷(1833年)，第296页。参见与我们的解释相应的第249节对警察的规定：“作为一种保护和保障大量特殊目的和利益的外部秩序与机构，警察的职能首先是实现和维持包含于市民社会的特殊性之中的普遍物，使这些特殊目的和利益存在于这种普遍物之中。其次，它作为上级领导，也要照顾超出这个社会之外的利益（第246节）。”基于黑格尔这里的前提，洛伦茨·冯·施泰因在1850年后写下了他的重要理论，并且为国家管理的社会进行辩护（《行政学》，1869年及以后）。}作为“普遍物的保证力量”（如果愿意的话，也可以叫作市民社会的“概念”的保证力量），警察要与市民社会的“不可阻挡的活动”、“快速发展的人口和工业”、“财富的积累”（第243节）以及“大量群众下降到一定生活水平之下”（第244节）作斗争。这里，浮现在黑格尔眼前的——首先是英国（第245节注释）——是“产生了贱民”的“群众”。在他看来，这些人将会炸裂作为“市民的”（亦即伦理上有序、政治上受到约束的）社会的社会形式。以前，野蛮的原始民族、家奴、日工和手工艺人处于旧的市民社会以外，现在则是贱民。但在黑格尔这里，与政治学的哲学传统不同，贱民不再表现出这种传统由以获得其历史合法性的积极限制；相反，贱民的存在要求国家行政（即“警察”）去限制社会的“不可阻挡的活动”，并用市民社会去整合他（因此，黑格尔的市民社会还不是1840年左右的早期社会主义者所言的固化了的阶级社会、“资产阶级社会”或者“布尔乔亚社会”）。

市民社会的第二种制度，即同业公会，既与市民社会的解放有关，也与政治蜕化为警察有关。如果说警察是古典政治的“降格”形态的话，那么同业公会就会让人想到古老的家庭团体(Oikos-Verband)，想到“家庭社会”。随着市民社会的内部转型，家庭的“经济上的”和伦理-实体性的秩序失去了基础，因此黑格尔也就不再像康德那样讨论“家庭社会的法”，而只谈到在“作为婚姻的它的直接概念的形态”，在其“外部定在、所有权和财产”，在“子女教育和家庭解体”之中的“家庭”（第160节以下）。黑格尔《法哲学原理》表达的是崭新的、不再与家庭经济细胞相联的家庭概念，并且意识到了为“经济学”的市民社会做好了准备的历史变化。因为有别于政治学的古典传统（对这种传统来说，市民社会居于家庭社会之上），在黑格尔看来，家庭在“市民社会中是一个从属的东西，只构成基础：它不再有如此广泛的活动范围。反之，市民社会则是将人扯到自己这里来的巨大力量，它要求人们为它劳动，人们的一切都通过它，并依靠它而活动”。\footnote{参见《全集》第8卷，第238节，补充，第299页。}在现代，家庭社会的成员从家庭中被“扯了出来”，成为“市民社会之子”，市民社会用它自己的土地“代替了个人赖以为生的父辈土地”，使整个家庭的存在都依赖于它，听任“偶然性”支配。\footnote{同上书，第298页：“首先，家庭是实体性整体，它的职责在于照顾个体的特殊方面。一方面为他提供手段和技能，使他能够从普遍财富中获利；另一方面，在他失去工作能力时，维持他的生活和赡养。但是，市民社会把个体从这种联系中揪出来，使家庭成员之间彼此生疏，并承认他们都是自立的个人。”}

现在，黑格尔认为市民社会的运动原则，它在经济学领域中的“广泛活动”，首先在“营利等级”、市民等级中得到实现。农村大量地保存了旧的经济形式、家庭生活的“实体性”，而在城市中，“个人没有等级尊严，通过它的孤立而沦为利己的工商业者，他的生活和享受也极不稳定”。\footnote{参见《全集》第8卷，第250节和第253节，以及与第256节国家的推演有关的社会的城乡划分：“前者是市民工商业、完全沉入自身并且个别化的反思的所在地，后者则是以自然为基础的伦理的所在地”。}在黑格尔看来，同业公会就是这种持存者，是在“市民社会的劳动制度”及其运行之中的唯一稳定者，因此就是那些劳动中孤立起来的个体的“共同体”。因此，这种稳定物就不是一个固定的等级。因为在这里，在市民社会的劳动制度中，“由于等级不再存在”，所以没有人还在按其等级而生活。\footnote{同上书，第253节，第309页。}毋宁说，同业公会是旧的家庭共同体的延伸。如第252节所言，它作为“第二家庭”出现，以反对社会运动的偶然性和特殊性，“对于一般的、比较远离个人及其特殊急需的市民社会来说，这种地位有待于进一步规定”。在它的这种地位中存在着其实存的“伦理根据”。尽管控诉“贱民的产生与之相联的工商业阶级的奢侈浪费”以及“劳动的日益机械化”（第253节），这种伦理根据似乎向黑格尔证明了它的正当性。

显然，同业公会和警察相对于市民社会的这种“地位”（采用黑格尔的说法）是与18世纪以来市民社会在现代世界中变化了的地位相吻合的。因为如开头想要表明的那样，旧市民社会从自身出发，包含了一种与经济和政治领域的关系，它的形态即产生于这些领域。旧的市民社会坐落在家庭社会及其归属的经济之上，并与之相对照，生长在政治即公共领域的土壤，只有“经济上”独立的社会成员才能进入这一领域；相反，黑格尔第一次在原则上将现代市民社会作为主题，并将其提高到概念意识中。此种市民社会将革命以前的旧欧洲世界的实体性的经济和政治秩序结构降格为偶性，其所以如此，是因为它自己成为了私人生活和公共生活的实体。社会将其政治部分交给国家，仅仅从传统政治那里保留了警察，作为它自己基本运动的管理和秩序功能；随着工业企业的发展，旧家政学的“稳定的东西”、家庭团体解体了，因此市民社会试图将那些在城市的劳动分工中孤立存在的个人重新联合为一个共同体。这样，随着革命而得到解放的社会便在警察和同业公会中恢复了传统的伦理—政治因素，基于此，它也就能够将自己建立为一个“市民社会”。我们可以跟黑格尔一起说，它们是“市民社会的无组织分子绕之旋转的”枢轴。\footnote{《全集》第8卷，第255节，第310页。在黑格尔这里，这句话仅涉及家庭伦理（婚姻的神圣性）和同业公会（个人在同业公会中的尊严），但同时一般而言，也为市民社会第三部分奠定了基础。}通过政治学传统加给其制度的限制，黑格尔阻断了市民社会成为“需要的体系”所勾勒的现代经济社会的趋势。因为社会作为需要体系就是无组织的市民社会，由此也就不再是黑格尔所理解的并在《法哲学原理》中加以阐述的市民社会的概念了。黑格尔通过在概念上将现代的决定性的历史政治过程——社会与国家的分离以及社会中的传统和现代的交互关系——表达出来，便能够思考市民社会的现代形态，同时通过更加古老的结构去限定其实质力量——尽管只是一个早已过去了的时代。可以说，唯有在这里，在一个既进行中介同时又运动着的、突破了中介的中心，法哲学的“市民社会”才历史性地产生出来，并且具有了其真正的概念地位。